A grant has to hold against the app that holds it and against whoever submits
its transactions: A1, the grantee's key, honest or stolen, and A2, a submitter
or relayer. Table~\ref{tab:tab-grantee} gives what each can do, what holds and
what remains. In Saltzer and Schroeder's terms~\cite{saltzer-schroeder},
\code{\_\_check\_auth} mediates every \code{require\_auth} on the account that
an authorization entry satisfies; one satisfied by invoker authority, as in the
owner's \code{execute}, never reaches it.

\perchtable{tab-grantee}

\paragraph{Goals and trust.} Containment (G1), permission soundness (G2) and
window safety (G3) are proved for lowering and the decision function, the code
around them tested; fail-closed refusal (G6) is tested; recovery invalidation
(G4) is proved in an account model (Theorem~10 of the verification companion),
tested and run on testnet. Other goals and adversaries are in the vision
companion. We trust the Soroban host and validators, the verifiers documents
name, the adopted recovery controller, the document compiler, rustc
and the wallet's display, and assume unforgeable signature schemes (Ed25519;
ECDSA P-256 for passkeys) and a collision-resistant SHA-256, under which the
verification companion proves on paper that an external signature binds the
account, the entry and the rule IDs. Budgets across grants, effects outside authorization,
confidentiality and verifier correctness are non-goals.
